% !TEX root = ../main.tex
\section{Predictions and a research agenda}\label{sec:agenda}
The argument yields hypotheses, not established treatment-selection rules.
\begin{enumerate}
\item \textbf{P1 (skill-level moderation).} Baseline executive function predicts acquisition of restructuring skill more strongly than acquisition of exposure-based skills, but predicts overall outcome only weakly.
\item \textbf{P2 (stress-dependence).} The gap between calm and in-the-moment reappraisal is larger in low-executive-function samples and at higher arousal than under the moderate stressors studied so far.
\item \textbf{P3 (component).} In adults with ADHD and anxiety, an exposure-and-behaviour protocol is non-inferior to the same protocol plus restructuring, under a clinically justified, prespecified margin.
\item \textbf{P4 (ADHD).} ADHD status predicts less benefit from the cognitive module specifically, rather than necessarily less benefit from the full CBT package.
\item \textbf{P5 (practice).} For persistent avoidance under high executive demand, structured behavioural practice produces more durable functional improvement than restructuring without comparable practice.
This prediction concerns a defined behavioural target, not all distress or every patient.
\end{enumerate}
A component comparison should include restructuring, behavioural practice and their combination, with comparable therapist contact and credible rationales.
Measure actual practice, expectancy change, skill acquisition, executive demand and functioning rather than inferring mechanisms from treatment labels.
Measure intellectual understanding, reported threat, avoidance and functioning separately, including change outside sessions.
Acceptance-related change should include valued action despite distress, not only reductions in distress intensity.
Greater emotional intensity during a session is not by itself evidence of durable change.
To test sequencing, hold cognitive content and contact as comparable as possible, distinguish early alliance from prior symptom change, and measure engagement alongside follow-up outcomes.
A final-phase cognitive consolidation hypothesis also requires outcomes beyond the final session, rather than treating a persuasive explanation as evidence of lasting change.
\textbf{What would weaken the hypotheses:} a precisely estimated absence of the predicted executive-function interaction would weaken P1.
A precisely estimated absence of a larger calm-versus-stressed reappraisal gap in low-executive-function samples would weaken P2.
A clinically meaningful restructuring benefit in the specified population would contradict P3.
A precisely estimated absence of a difference in cognitive-module benefit by ADHD status would weaken P4.
A precisely estimated absence of a practice advantage would weaken P5.
A non-significant difference alone is insufficient: the relevant superiority or non-inferiority contrast, margin, outcome, follow-up and power must be specified before data collection.
Mediation analyses also require temporal ordering and assumptions about confounding; correlated process change is not enough.
The immediate research priorities are a structured umbrella search with independent screening, analyses of existing participant-level data where available, and adequately powered component/moderator designs.
Any prospective clinical comparison would require ethical review and appropriate treatment provision; none is conducted here.

\section{Conclusion}
The ten reasons graded in Table~\ref{tab:reasons} do not support the claim that CBT generally fails.
They identify substantial non-response, relapse, under-recorded harms, comparator-sensitive effects, uncertain mechanisms, variable delivery and unequal access.
Most of these limitations are shared with other psychological treatments.
The executive-function argument remains narrower than its popular formulation: reappraisal uses working memory, but stress does not reliably abolish it, and clinical moderation remains insufficiently tested.
Component evidence does not establish that explicit restructuring is always necessary, always dispensable, or the unique route to cognitive change.
Evolutionary conservation motivates the learning-based hypothesis, whose clinical value requires direct testing.
Acceptance, alliance, reactance and feedback research further caution against treating treatment content independently of goals and delivery.
Researcher allegiance and publication selection warrant scrutiny, but do not establish CBT-specific reviewer suppression or invalidate the efficacy literature.
The useful next question is which components, delivered in which way, help which patients achieve durable functional change.
